\documentclass[a4paper]{scrartcl}
\usepackage[utf8]{inputenc}
\usepackage{amsmath,amssymb,amsthm}
\usepackage{enumitem,setspace}
\usepackage[usenames,dvipsnames,table]{xcolor}
\usepackage{graphicx}
\DeclareGraphicsExtensions{.pdf,.png,.eps,.jpg}
\usepackage{algorithm2e}

% --- math operators ---
\usepackage{mathtools}
\DeclarePairedDelimiter\set{\{}{\}} % use $\set{1, 2, 3}$
\DeclarePairedDelimiter\abs{\lvert}{\rvert}
\DeclarePairedDelimiter\norm{\lVert}{\rVert}
\DeclarePairedDelimiter\ceils{\lceil}{\rceil}
\DeclarePairedDelimiter\floor{\lfloor}{\rfloor}
\def\then{\ensuremath{\Rightarrow}} % =>
\def\iff{\ensuremath{\Leftrightarrow}} % if and only if <=>
\def\to{\ensuremath{\rightarrow}} % ->
\def\ot{\ensuremath{\leftarrow}} % ->
\def\Oh{\ensuremath{\mathcal{O}}} % big O like $\Oh(n)$
% ---
\DeclareMathOperator{\preorder}{preorder}
\DeclareMathOperator{\postorder}{postorder}
\DeclareMathOperator{\depth}{depth}

% --- FILL OUT THESE 3 FIELDS ---
\def\sheetNumber{1}
\def\names{NAMEN} 
\def\sumPoints{20} 

% --- ### BEGIN DOCUMENT ### ---
\begin{document} 

\begin{doublespace} % header
\noindent \textbf{\LARGE{Übungsblatt \sheetNumber}}
\hfill  {\large Punkte: $\boxed{\qquad  /\; \sumPoints}$}\\
{\Large \textbf{Fortgeschrittene Algorithmen (Wintersemester 2022/23)}\\
Bearbeitet von: \textbf{\names}\\
\today}
\end{doublespace}


\section*{Aufgabe 1}

\section*{Aufgabe 2}

\section*{Aufgabe 3}
a) Der algorithmus terminiert garantiert da in jedem Schritt die Menge S um ein Element verkleinert wird bis \\
nur noch ein Element in S ist. Der Algorithmus ist korrekt.\\

b) Jeder Rekursive Aufruf betrachtet eine Teilmenge [S, S-1 S-2...1]. Im worst case werden alle \\
Elemente betrachtet also n Rekursive Aufrufe. In jedem Aufruf wird eine Konstante Zeit benötigt. \\
D.h. $T (n) = T (n-1) + (n - 1) \to T (n) = \Oh (n^2)$ \\

c) Nur wenn das kleinste Element ausgewählt wird, wird die überprüfung in Schritt 5 des Algorithmus \\
durchgeführt. Dies passiert mit einer Wahrscheinlichkeit von $\frac{1}{n}$ im ersten Aufruf. \\
In diesem Fll liegt die Laufzeit bei $\Oh (n)$. In allen anderen Fällen wird \\
wird nur der Rekursive Aufruf durchgeführt. Dies passiert mit einer Wahrscheinlichkeit von \\
$\frac{n-1}{n}$. \\
Die erwarte Laufzeit $T (n)$ lässt sich also wie folgt beschreiben: \\
$T(n) = T(n-1) + \frac{1}{n} \cdot \Oh (n)$ \\
$T(n) = T(n-1) + \Oh (1)$ \\
$T(n) = \Oh (n)$ \\
\section*{Aufgabe 4}


% -- code for figures
% \begin{figure}[htb]
% \centering
% \includegraphics{filename} % use option [width=\linewidth] to change size to linewidth (or 0.6\linewidth for example)
% \label{fig:name}
% \end{figure}


\end{document}
% --- ### END DOCUMENT ### ---